\begin{trapbox}
\textbf{易错点1: 忘记两个斜率必须都存在}

公式 $k_lk_{OP}=-b^2/a^2$ 的使用前提是 $x_0y_0\ne0$.
\begin{itemize}
  \item $P=(\pm a,0)$: $k_{OP}=0$, 切线 $x=\pm a$ 的斜率不存在.
  \item $P=(0,\pm b)$: $OP$ 斜率不存在, 切线 $y=\pm b$ 的斜率为 $0$.
\end{itemize}
不可把“$0\times\infty$”当作 $-b^2/a^2$.
从非轴点趋向顶点时\textbf{乘积的极限}虽为 $-b^2/a^2$,
但顶点处\textbf{乘积本身没有定义}.

\textbf{易错点2: 误认为 $OP$ 就是椭圆的法线}

切线的一个法向量为
$\left(x_0/a^2,y_0/b^2\right)$, 而 $OP$ 的方向向量为 $(x_0,y_0)$.
当 $a\ne b$ 且 $x_0y_0\ne0$ 时, 二者\textbf{不平行}.
不要把圆“半径垂直于切线”的性质直接搬到一般椭圆上.

\textbf{易错点3: 错用垂直判定 $k_1k_2=-1$}

只有\textbf{两条斜率都存在}时, 才能以斜率积 $-1$ 判断垂直.
非轴点处椭圆切线与 $OP$ 的斜率积为 $-b^2/a^2$,
由于 $a>b>0$, 有 $-1<-b^2/a^2<0$, 故不垂直.
四个轴顶点处却仍然垂直, 因为一条直线水平、另一条竖直.

\textbf{易错点4: 把 $e^2-1$ 记成 $1-e^2$}

椭圆中 $c^2=a^2-b^2$, 因而
\[
e^2-1=\frac{c^2-a^2}{a^2}=-\frac{b^2}{a^2}<0.
\]
正确关系是 $k_lk_{OP}=e^2-1$, \textbf{负号不能漏}.

\textbf{易错点5: 分不清长轴方向}

本题长轴沿 $x$ 轴, 对于 $x^2/a^2+y^2/b^2=1$,
结果为 $-b^2/a^2$.
若写成 $x^2/b^2+y^2/a^2=1$ ($a>b>0$),
同样方法给出 $k_lk_{OP}=-a^2/b^2$, 不能机械抄成 $-b^2/a^2$.
更通用的记法是: 对 $x^2/A+y^2/B=1$ ($A,B>0$),
当两个斜率存在时, 乘积等于 $-B/A$.

\textbf{易错点6: 中心平移后仍使用原点连线}

若椭圆为
$\frac{(x-h)^2}{a^2}+\frac{(y-k)^2}{b^2}=1$,
则应使用\textbf{中心 $S(h,k)$ 到切点 $P$} 的斜率
$k_{SP}=(y_0-k)/(x_0-h)$, 从而 $k_lk_{SP}=-b^2/a^2$.
一般不能把 $k_{SP}$ 换成 $k_{OP}$.

\textbf{易错点7: 忽略旋转坐标轴的影响}

以上简单定值适用于椭圆主轴\textbf{与所用坐标轴平行}的情形.
将椭圆旋转成带 $xy$ 交叉项的一般方程后,
沿原坐标轴计算的两条斜率通常不再满足同一个 $-b^2/a^2$ 定值;
需要先旋转回主轴坐标系, 或重新从切线方程推导.

\textbf{易错点8: 中点弦公式不等于任意一条弦都能直接代入}

弦 $AB$ 中点为 $M$, 在 $k_{AB}$、$k_{OM}$ 都存在时才可以写
$k_{AB}k_{OM}=-b^2/a^2$. 若弦经过椭圆中心, 则 $M=O$,
直线 $OM$ 无法定义; 若涉及水平弦或竖直弦, 也要改用方向判断.

\textbf{易错点9: 忘记“切点在椭圆上”}

公式 $\frac{x_0x}{a^2}+\frac{y_0y}{b^2}=1$ 作为\textbf{在 $P$ 处的切线方程}
要求 $P$ 在椭圆上. 若 $P$ 为椭圆外一点,
这个式子通常表示两条切线的\textbf{切点弦所在直线}（在实切线存在时）, 不是“过 $P$ 的切线”.

\textbf{易错点10: 只知道切线斜率就声称切线方程唯一}

同一斜率的非竖直椭圆切线通常有\textbf{两条}平行直线;
若求直线方程, 还需确定具体切点或截距.

\textbf{快速自检}: 在 $x^2/4+y^2=1$ 上取
$P=(\sqrt2,1/\sqrt2)$, 可得 $k_{OP}=1/2$, $k_l=-1/2$,
相乘为 $-1/4$, 符合 $-b^2/a^2$.
\end{trapbox}
